\begin{ouverture}[calcul]{4}{La division : entière et décimale}{Partie 1 · Nombres entiers et calculs}
{Quotient · Reste · Quotient décimal exact}
\objectif{Poser une division euclidienne et contrôler son résultat.}
\objectif{Poursuivre une division après la virgule.}
\objectif{Choisir entre quotient entier et quotient décimal selon le contexte.}
\prerequis{Division euclidienne du CM1 et rangs décimaux.}
\end{ouverture}
\pagestyle{otMain}\essentiel\fontsize{11}{15}\selectfont
\phantomsection\label{cours-4-1}
\essSection{1}{La division euclidienne, bien installée}
\begin{definition}Diviser 1\,258 par 6, c’est écrire $1\,258=6\times209+4$ : quotient 209, reste 4, avec un reste toujours plus petit que le diviseur. Le contrôle refait le chemin inverse : $6\times209+4=1\,258$.\end{definition}
\begin{savoirfaire}[Les zéros du quotient]Poser proprement : chercher combien de fois le diviseur entre dans les premiers chiffres, abaisser chiffre à chiffre, placer un 0 au quotient quand le chiffre abaissé ne suffit pas : $3\,618:9=402$.\end{savoirfaire}
\begin{exemple}Selon la question, le même calcul répond différemment : avec 230 œufs et des boîtes de 6, on remplit 38 boîtes ($230=6\times38+2$), mais il faut 39 boîtes pour tout transporter.\end{exemple}
\begin{attention}Un reste se compare toujours au diviseur : si le reste le dépasse, un paquet de plus est possible et la division n’est pas terminée.\end{attention}
\clearpage
\phantomsection\label{cours-4-2}
\essSection{2}{Continuer après la virgule}
\begin{definition}Quand le partage doit être exact, la division se poursuit dans les rangs décimaux : $27:4$ donne quotient 6 et reste 3 ; les 3 unités deviennent 30 dixièmes ; $30:4$ donne 7 dixièmes, reste 2 ; les 2 dixièmes deviennent 20 centièmes ; $20:4=5$. Donc $27:4=6{,}75$ exactement.\end{definition}
\begin{exemple}Avec un dividende décimal, même méthode : $8{,}7:6$. Les 8 unités donnent 1 ; reste 2, soit 27 dixièmes avec le 7 ; $27:6=4$ reste 3, soit 30 centièmes ; $30:6=5$. Donc $8{,}7:6=1{,}45$. La virgule du quotient se place en passant aux dixièmes.\end{exemple}
\begin{savoirfaire}[Contrôler le quotient]Contrôler un quotient décimal par la multiplication : $6{,}75\times4=27$ et $1{,}45\times6=8{,}7$. Estimer aussi : $27:4$ est entre 6 et 7.\end{savoirfaire}
\begin{attention}Certains partages ne s’arrêtent jamais : $10:3=3{,}333\ldots$ On s’arrête alors au rang demandé par l’énoncé. Un quotient décimal exact existe quand le reste finit par tomber à zéro.\end{attention}

\clearpage\rubrique{Découverte}
\begin{blocLivre}{ex-4-1}\begin{exercice}[1]{Une division posée}
Poser la division euclidienne de $1\,258$ par 6, écrire l’égalité et la contrôler.
\end{exercice}\end{blocLivre}
\begin{blocLivre}{ex-4-2}\begin{exercice}[1]{Le zéro au quotient}
Calculer $3\,618:9$, puis $2\,814:7$.
\end{exercice}\end{blocLivre}
\begin{blocLivre}{ex-4-3}\begin{exercice}[1]{Après la virgule}
Calculer le quotient décimal exact de $27:4$.
\end{exercice}\end{blocLivre}
\begin{blocLivre}{ex-4-4}\begin{exercice}[1]{Dividende décimal}
Calculer $8{,}7:6$.
\end{exercice}\end{blocLivre}
\begin{blocLivre}{ex-4-5}\begin{exercice}[1]{Moitiés et quarts décimaux}
Calculer $9:2$, puis $9:4$.
\end{exercice}\end{blocLivre}

\rubrique{Entraînement}
\begin{blocLivre}{ex-4-6}\begin{exercice}[2]{Choisir le bon quotient}
On partage 75 euros entre 4 personnes. On range 75 cahiers dans des cartons de 4. Donner la réponse adaptée à chaque situation.
\end{exercice}\end{blocLivre}
\begin{blocLivre}{ex-4-7}\begin{exercice}[2]{Le reste et l’égalité}
Sans poser la division, dire si cette égalité décrit une division euclidienne : $507=8\times62+11$. Corriger si besoin.
\end{exercice}\end{blocLivre}
\begin{blocLivre}{ex-4-8}\begin{exercice}[2]{Division par dix, cent, mille}
Calculer $4\,500:100$, puis $38:10$, puis $7\,200:1\,000$.
\end{exercice}\end{blocLivre}
\begin{blocLivre}{ex-4-9}\begin{exercice}[2]{Encadrer le quotient}
Sans calculer le quotient exact, encadrer $250:7$ entre deux entiers consécutifs, puis calculer le quotient euclidien et le reste.
\end{exercice}\end{blocLivre}
\begin{blocLivre}{ex-4-10}\begin{exercice}[2]{L’erreur de virgule}
Sami calcule $14{,}4:6=2{,}04$. Vérifier par la multiplication, expliquer l’erreur, corriger.
\end{exercice}\end{blocLivre}

\rubrique{Problèmes}
\begin{blocLivre}{ex-4-11}\begin{exercice}[3]{La course en relais}
Une équipe doit parcourir 15 km en 4 relais de même longueur. Quelle distance chaque coureur parcourt-il ?
\end{exercice}\end{blocLivre}
\begin{blocLivre}{ex-4-12}\begin{exercice}[3]{Les packs d’eau}
Un magasin reçoit 1\,350 bouteilles à ranger par packs de 8.\begin{questions}
\item Combien de packs complets prépare-t-on ?
\item Combien de bouteilles restent seules ?
\item Combien manque-t-il de bouteilles pour un pack de plus ?
\end{questions}
\end{exercice}\end{blocLivre}
\begin{blocLivre}{ex-4-13}\begin{exercice}[3]{Le ruban partagé}
Un ruban de $9{,}6$ m est coupé en 8 morceaux identiques. Quelle est la longueur d’un morceau, en mètres puis en centimètres ?
\end{exercice}\end{blocLivre}
\begin{blocLivre}{ex-4-14}\begin{exercice}[3]{La cagnotte des quatre}
Quatre amis se partagent équitablement une cagnotte de 58 euros, en utilisant aussi les centimes. Combien reçoit chacun ? Vérifier que le partage rend bien toute la cagnotte.
\end{exercice}\end{blocLivre}
\begin{blocLivre}{ex-4-15}\begin{exercice}[3]{Un partage impossible ?}
Lila veut partager exactement 10 mètres de ficelle en 3 morceaux de même longueur, en écrivant la mesure en mètres. Que se passe-t-il ? Proposer une réponse raisonnable au centimètre près.
\end{exercice}\end{blocLivre}

\clearpage
